\scp{Methodological issues relating to marsupial terms}{14-02}
In the comparative method,
innovations that are shared in two or more lower level subgroups
can be ascribed to the common parent language.
So if daughter group X has innovations that are also found in
daughter group Y, but not found elsewhere,
then the reconstruction that accounts for the forms in both groups
can be assigned to their common parent language,
or the node that encompasses both daughter groups.
\citet{Blust1982} found two etyma that in his view were shared
exclusively between CMP and EMP, but not found in WMP languages.
He reconstructed those as **kandoRa and **mans(aə)r.
Given that marsupials are not found west or north of Sulawesi,
it would have been unreasonable to assign these forms to PMP,
which includes languages of the Philippines,
western Indonesia and Malaysia where there are no marsupials.
Therefore, he concluded, they must belong to an intermediate
node that falls within the distribution of marsupial terms.
There are many reasons to not force the languages of Sulawesi
to be a subgroup with Proto-Oceanic,
and Blust found no cognates in Sulawesi.
So given Blust’s assumptions about subgrouping,
the solution was to propose an intermediate node (PCEMP) that
included CMP and EMP, and excluded WMP, to which Sulawesi languages were assumed
to belong (see \frf{01-02}).

We also noted in \srf{02-04-04} that \citet{Schapper2011} had
challenged the assumptions, interpretations of the data,
and subgrouping conclusions of \citet{Blust1982}.
She pointed out errors in glossing,
the highly restricted distribution of reflexes of **kandoRa,
and argued that **kandoRa should be limited to EMP.
She pointed out that cognates of **mans(aə)r in most languages
outside of EMP actually mean ‘cuscus’,
not ‘bandicoot’, and consequently suggested that PCEMP **mans(aə)r
should be glossed ‘cuscus’,
later shifting to ‘bandicoot’ for Oceanic languages.
In spite of conceding that Schapper was correct on many of her
criticisms and corrections, \citet{Blust2012b} nevertheless
maintained his (\citeyear{Blust1982}) position and glosses.

While discussing Austronesian incursions into the region,
in \srf{02-01} we briefly noted \citeauthor{Blust2012b}'s (\citeyear[74]{Blust2012b}) rejection
of Schapper's “implicit claim” that both the glosses in the languages
and the limited distribution of bandicoots show the Austronesian-speaking
peoples encountered the marsupial cuscus before they encountered
the bandicoot, saying her scenario “implies an entry via Sulawesi into the Lesser
Sundas, with subsequent eastward movement into the Moluccas”,
for which he said there was no linguistic evidence.
In \crf{ch:SB}, we presented compelling evidence in the historical
phonology for at least one incursion from Sulawesi eastward into
Maluku, from central Sulawesi through Banggai to Taliabu.\footnote{
		The distribution of the \it{Strigocuscus pelengensis} ‘Banggai
		(dwarf) cuscus’, found in the Banggai Islands
		and the nearby Sula Islands,
		parallels the linguistic link between \tsc{Taliabu}
		and Banggai discussed in \crf{ch:SB}.
		The distance between Sanana
		and Buru appears to have been too great for that species to cross.}
 Evidence from the lexicon of additional Sulawesi-Maluku links
are presented in this chapter.

Surveying methodological problems when dealing with old sources
in the region, in \srf{02-04-04}, we briefly touched on the care needed when
filtering glosses through various languages.
We noted that \citet{Schapper2011} pointed out that several of
the languages that \citet{Blust1982} had glossed as ‘bandicoot’,
were actually outside the region where bandicoots are found.
And we provided independent evidence that the Dutch gloss \it{buidelrat}
(literally ‘pouch rat’) used in his sources was equated directly
in Dutch sources of that era with Malay terms for ‘cuscus’,
and not ‘bandicoot’, as Blust had thought.

Given the extremely limited distribution of reflexes of **kandoRa
near New Guinea, and that almost all cognates of **mans(aə)r in Maluku actually
mean ‘cuscus’ and not ‘bandicoot’ as \citet{Schapper2011} correctly pointed
out, to understand \citeauthor{Blust2012b}'s (\citeyear{Blust2012b}) insistence on maintaining **kandoRa
for all of CEMP,
and resistance to reglossing **mans(aə)r from ‘bandicoot’ to
‘cuscus’, one must understand his subgrouping hypothesis
and assumptions about the breakup
and dispersal of PCEMP into CMP and EMP.
In Blust's view, the breakup of PCEMP happened somewhere in “northeast Indonesia”.
The group that eventually became \tsc{POc} travelled eastward
and dispersed from there throughout Oceania.
\citet[98]{Blust2008} says, “As noted in \citet{Blust1993},
Central Malayo-Polynesian almost certainly was the product of
rapid movement from northeast Indonesia through the central
and southern Moluccas and into the Lesser Sunda islands,
forming a dialect chain that favoured the ready diffusion of
innovations across incipient language boundaries,
and hence destroying any tree-like structure that may have previously
arisen.” 
Given that back-eddy type of scenario dispersing from ``northeast
Indonesia'' that excludes Sulawesi,
in Blust's perspective the problem of cognates of **kandoRa only
being found in a very small number of languages close to the
EMP boundary is trivial.
That is where they started,
and they just did not move very far south or west.
If one assumes some of the Austronesians came from the west or
northwest moving eastward, then Blust's position is confusing.
As yet there is little or no linguistic or archaeological evidence
to support Blust’s purported migration scenario from northeast
Indonesia travelling south and west through central Maluku,
southern Maluku, Lesser Sundas.
(See \srf{15-06-03}, and \citealp[4]{Bellwood2019b}.)

It turns out that the stories of both **kandoRa
and **mans(aə)r are far more complex
and interesting than those presented in \citet{Blust1982,Blust2012b},
and even more far-reaching than the story told by \citet{Schapper2011}.
The stories of both of these terms link to various Papuan language families.
And the stories of both of these terms have ties to Sulawesi,
which lies within the biogeographical region where marsupials
are naturally found, but outside the region where bandicoots are found,
and outside the purported “CEMP” region.
Reviewing the data and the issues, and going beyond those works
have significant implications for both methodology
and subgrouping, and thus warrants doing so in detail.
That is the purpose of this chapter.
A detailed list of terms for marsupials
and related animals is found in appendix \ref{ch:SurMar},
which has 484 terms for marsupials
(and similar mammals) from 284 languages/varieties,
of which 227 are Austronesian,
and 57 are Papuan.

There are two families of words within the region: \it{k-} initial
(Blust’s **kandoRa, our **kVndoR(a)) and \it{m-} initial (Blust’s **mans(sə)r, our **mantəR).
The various words that can be assigned to each of these families
are incrementally similar, regardless of subgroup or language family,
and contact has almost certainly taken place to lead to their modern distribution.
We are not proposing that these are reconstructions from a single (Austronesian) parent language.
We note there are other reconstructable forms not in focus here,
particularly in the Sulawesi region.
The discussion and data in the following two sections show that
within our target region:

\newpage
\begin{enumerate}
	\item	Both **kVndoR(a)
				and **mantəR should be glossed ‘cuscus’.
	\item	There is only one language outside of Oceanic languages
				having the gloss ‘bandicoot’ for a possible,
				but anomalous, reflex of **mantəR.
	\item	There is only one language outside of EMP reflecting all
				three syllables and vowels of **kandoRa.
	\item	Languages are in a kind of complementary distribution. They either
				have a term from the \it{k-}initial family of forms (e.g.
				**kVndoR(a)), or the \it{m-}initial family (**mantəR),
				but not both (explained further below)
	\item	Languages reflecting **kVndoR(a) are extremely limited
				in their distribution, but link incrementally to formally similar words in Sulawesi
				to the west, and to Papuan languages in the north and east.
	\item	Languages reflecting **kVndoR(a) are not found in most Wallacean subgroups.
				They are only found in one node of one subgroup,
				and within that node there is great variation.
	\item	The form **mantəR (with close variations from subgroup
				to subgroup) better accounts for the sound correspondences across
				the region than does **mans(aə)r.
	\item	Reflexes of **mantəR ‘cuscus’ are also found in a number
				of languages in Sulawesi,
				outside the purported CEMP region.
\end{enumerate}

The existence of a subgroup,
particularly one which has featured importantly in understanding
the spread and dispersal of Austronesian languages,
should not be founded on such limited
and marginal lexical data.

As mentioned above, we note that reflexes of the two types of
forms (those beginning with initial \it{k-},
and **mantəR ‘cuscus’) seem to be in complementary distribution.
We find several types of complementary distribution:

\begin{enumerate}
	\item	Except for one Wallacean subgroup,
				reflexes of **kVndoR(a) and **mantəR are not found together.
				Within that one subgroup,
				reflexes of **kVndoR(a) are found in only one node,
				and reflexes of **mantəR are not found in that same node.
				Reflexes of **kVndoR(a) are not found in any other Wallacean subgroup.
	\item	Words related to **kusay in Sulawesi refer to ‘bear cuscus’.
				Sulawesi languages have a separate term,
				**taŋali, for the ‘dwarf cuscus’.
				Words related to **mantəR in \tsc{Muna-Buton} (\tsc{Celebic}) languages
				refer to ‘bear cuscus’.
				In Wallacean subgroups **mantəR is a generic term for ‘cuscus’.
	\item	The few languages that may appear to have both the \it{k}-initial
				family of forms and a reflex of **mantəR have them in different dialects.
				For example, most varieties of Cia-Cia in Southeast Sulawesi
				have only \it{mante} ‘bear cuscus’ (from **mantəR),
				but the Wasambua dialect of Cia-Cia has only \it{kuse} ‘bear cuscus’ (from **kusay).
\end{enumerate}